\chapter*{Abstract}
\addcontentsline{toc}{chapter}{Abstract}
\markboth{Abstract}{Abstract}

Financial markets do not behave homogeneously over time. Instead, they
alternate between periods of sustained growth, prolonged downturns, and
episodes of extreme instability in which volatility spikes and
cross-asset correlations break down. These ``market regimes'' ---
informally described as \emph{bull}, \emph{bear} and \emph{crisis}
markets --- have a direct impact on risk management, asset allocation
and investment decision-making. Identifying, in an objective and
automated manner, which regime the market is currently in is therefore a
problem of substantial academic and practical interest.

Traditional approaches to this problem rely mainly on price-based and
technical variables --- returns, volatility, moving averages, momentum
--- combined with statistical methods such as regime-switching models or
unsupervised learning techniques like \textit{K-Means}. However, these
variables are, by construction, \emph{backward-looking}: they reflect
price movements that have already occurred. At the same time, the rise of
\textit{Transformer}-based language models, and in particular of
finance-specific models such as \textit{FinBERT}, has made it possible to
systematically extract the emotional tone --- positive, negative or
neutral --- of large volumes of financial news, providing a
\emph{contemporaneous} source of information that is potentially
complementary to classical technical indicators.

This Bachelor's Thesis designs and implements a complete end-to-end
pipeline that combines both sources of information --- price/technical
indicators and sentiment extracted from financial news --- with the goal
of identifying market regimes through unsupervised clustering techniques,
and of rigorously assessing whether incorporating sentiment improves the
quality and economic interpretability of the resulting regimes compared
to a purely technical approach.

\section*{Objectives}

The overall objective of this project is to \textbf{develop and validate
a clustering-based market regime detection methodology, comparing a model
that relies exclusively on S\&P 500 technical indicators against models
that additionally incorporate sector-level sentiment extracted from
financial news using FinBERT}. This overall objective is broken down into
the following specific objectives:

\begin{itemize}
  \item Build a unified dataset combining daily technical indicators for
  the S\&P 500 index with a sector-level (GICS) sentiment indicator
  derived from the \textit{FNSPID} financial news corpus.
  \item Design and implement a reproducible sentiment-scoring procedure
  based on the \textit{FinBERT} model, together with a sector-level
  aggregation and temporal smoothing scheme.
  \item Build three alternative \textit{K-Means} clustering models ---
  technical-only, technical+sentiment and smoothed-sentiment-only --- and
  compare them using internal validation metrics (silhouette coefficient,
  principal component analysis) and feature-importance analysis.
  \item Incorporate a \textit{Hidden Markov Model} (HMM) as a secondary
  comparison method, evaluating its differences with respect to
  \textit{K-Means} in terms of temporal stability and economic
  differentiation of the regimes.
  \item Economically validate the obtained regimes by computing
  annualized return, volatility, Sharpe ratio, drawdown and other
  metrics, broken down by GICS sector and by regime.
\end{itemize}

\section*{Methodology}

The developed pipeline, schematically represented in
\autoref{fig:abstract-pipeline}, consists of the following main stages:

\begin{enumerate}
  \item \textbf{Data acquisition and filtering}: the starting point is
  the public \textit{FNSPID} (Financial News and Stock Price Integration
  Dataset) corpus, which contains financial news headlines associated
  with US-market tickers. News items are filtered to retain only those
  associated with companies that were constituents of the S\&P 500 at the
  time of publication, and are reassigned to the next open market day to
  avoid look-ahead bias.
  \item \textbf{Sentiment scoring with FinBERT}: each headline is
  processed with the \textit{ProsusAI/FinBERT} model, yielding
  probabilities $p_{\text{pos}}$, $p_{\text{neg}}$ and $p_{\text{neu}}$. A
  continuous sentiment score $s = p_{\text{pos}} - p_{\text{neg}} \in
  [-1,1]$ is defined, together with a categorical label
  (positive/negative/neutral) obtained as the argmax of the three
  probabilities.
  \item \textbf{Sector-level aggregation}: each ticker is mapped to its
  corresponding GICS sector (11 sectors), and for each (date, sector)
  combination the mean and standard deviation of sentiment across all
  published news items are computed. Given the noisy nature of the daily
  signal, a 21-session moving average (\textit{MA21}) is additionally
  applied, which substantially reduces high-frequency noise while
  preserving medium-term trends in sector sentiment.
  \item \textbf{Technical indicators}: nine technical indicators are
  computed from daily S\&P 500 prices --- RSI, 50- and 200-session moving
  average ratios, 21- and 126-session momentum, 30- and 252-session
  realized volatility, and 30- and 252-session return skewness.
  \item \textbf{Clustering dataset construction}: technical indicators and
  sector-level sentiment variables (raw and smoothed) are merged into a
  single daily dataset, resulting in 1,132 observations spanning May 2009
  to February 2024.
  \item \textbf{K-Means clustering}: three models are trained ---
  \textbf{Model A} (technical indicators only), \textbf{Model B}
  (technical indicators + raw sector sentiment) and \textbf{Model C}
  (smoothed sector sentiment only) --- with $K=3$, selected via the elbow
  method and the silhouette coefficient. In addition, a \textbf{Gaussian
  HMM} (full covariance, $K=3$ states) is trained on the same feature sets
  as a comparison method.
  \item \textbf{Regime labeling and interpretation}: the resulting
  clusters are labeled as \emph{Bull}, \emph{Bear} or \emph{Crisis} based
  on centroid analysis (mean momentum, volatility and RSI of each
  cluster).
  \item \textbf{Economic validation}: for each combination of GICS sector
  and regime, annualized return, annualized volatility, Sharpe ratio
  ($r_f=2\%$), maximum drawdown, percentage of positive days and extreme
  daily returns (best and worst day) are computed.
\end{enumerate}

\begin{figure}[H]
  \centering
  % TODO: insert full pipeline diagram (block diagram)
  \fbox{\parbox{0.85\textwidth}{\centering\vspace{3cm}
  Block diagram of the complete pipeline: FNSPID
  $\rightarrow$ S\&P 500 filtering $\rightarrow$ FinBERT $\rightarrow$
  GICS sector aggregation (mean/std + MA21) $\rightarrow$ S\&P 500
  technical indicators $\rightarrow$ clustering dataset $\rightarrow$
  K-Means (Models A/B/C) and HMM $\rightarrow$ economic validation by
  sector and regime.\vspace{3cm}}}
  \caption{General overview of the pipeline developed in this project.}
  \label{fig:abstract-pipeline}
\end{figure}

\section*{Main results}

Silhouette analysis in a common principal-component space (retaining 80\%
of explained variance) shows that \textbf{Model A (technical)} achieves
the best internal cohesion (silhouette $\approx 0.193$), followed by
\textbf{Model B (technical + sentiment)} ($\approx 0.178$) and
\textbf{Model C (smoothed sentiment only)} ($\approx 0.162$). This
indicates that technical variables remain the strongest structural
discriminators of the data, and that sentiment alone is not sufficient to
define well-separated clusters.

However, \textbf{economic validation} reveals a more nuanced picture.
Model B, which incorporates sector sentiment alongside technical
indicators, produces regimes with \emph{sharper} economic differentiation
than Model A: the \emph{Bear} regime of Model B exhibits negative Sharpe
ratios across virtually all sectors (e.g., $-0.36$ for \textit{Financial
Services} and $-0.26$ for \textit{Energy}), while the \emph{Crisis} regime
counter-intuitively concentrates the highest annualized returns (up to
$+91.9\%$ for \textit{Real Estate}), reflecting high-volatility episodes
with strong rebounds. The \textit{Energy} sector in the \emph{Bear} regime
exhibits the most severe maximum drawdown of the entire analysis
($-74.9\%$), as well as the best and worst single-day returns of the whole
sample ($+16.9\%$ and $-24.7\%$, respectively), highlighting the extremely
volatile nature of this regime.

\textbf{Model C}, based exclusively on smoothed sentiment, produces
centroids that are nearly indistinguishable in technical variables (e.g.,
the 50-session moving-average ratio ranges from 1.0086 to 1.0134 across
regimes), constituting a \emph{meaningful negative result}: aggregated
sentiment alone, without price information, fails to capture market
regime dynamics with the same sharpness as technical indicators.

The comparison with the \textbf{HMM} model shows that, while both methods
identify qualitatively similar regimes, \textit{K-Means} produces a
\emph{much sharper} economic differentiation between regimes (Sharpe
ratios ranging approximately from $-3.1$ to $+3.8$), whereas the HMM
generates states with much milder Sharpe ratios (roughly between $-0.6$
and $+1.5$), consistent with its nature as a generative model with smooth
transitions between states.

\section*{Conclusions}

The results obtained confirm that technical indicators provide the
strongest foundation for clustering-based market regime detection, but
that incorporating sector-level sentiment extracted via FinBERT
\textbf{adds incremental value in terms of the economic differentiation}
of the resulting regimes, particularly for distinguishing sustained
\emph{bear-market} episodes from high-volatility \emph{crisis} episodes
with sharp rebounds. These findings open avenues for future work focused
on the dynamic combination of both information sources, the extension of
the analysis to other indices and markets, and the exploration of
sequential deep-learning architectures for regime detection.

\vspace{1cm}
\noindent\textbf{Keywords:} market regimes, clustering, K-Means, Hidden
Markov Model, sentiment analysis, FinBERT, financial news, FNSPID, S\&P
500, economic validation.

\clearpage
